\documentclass[11pt]{article}

% ------------------------------------------------------------
% Packages
% ------------------------------------------------------------
\usepackage[margin=1in]{geometry}
\usepackage{amsmath, amssymb, amsfonts}
\usepackage{enumitem}
\usepackage{booktabs}
\usepackage{longtable}
\usepackage{array}
\usepackage{xcolor}
\usepackage{hyperref}
\usepackage{titlesec}
\usepackage{fancyhdr}

% ------------------------------------------------------------
% Formatting
% ------------------------------------------------------------
\setlength{\parindent}{0pt}
\setlength{\parskip}{6pt}

\titleformat{\section}
  {\large\bfseries}
  {\thesection}
  {0.5em}
  {}

\titleformat{\subsection}
  {\normalsize\bfseries}
  {\thesubsection}
  {0.5em}
  {}

\titleformat{\subsubsection}
  {\normalsize\itshape\bfseries}
  {\thesubsubsection}
  {0.5em}
  {}

\pagestyle{fancy}
\fancyhf{}
\rhead{AS.110.106 Calculus I}
\lhead{Johns Hopkins University}
\cfoot{\thepage}

% ------------------------------------------------------------
% Custom commands
% ------------------------------------------------------------
\newcommand{\topic}[1]{\textbf{#1}}
\newcommand{\important}[1]{\textbf{Important:} #1}

% ------------------------------------------------------------
% Document
% ------------------------------------------------------------
\begin{document}

\begin{center}
    {\LARGE \textbf{AS.110.106 Calculus I}}\\[4pt]
    {\Large Biology and Social Sciences}\\[8pt]
    {\large Full Course Structure and Learning Framework}\\[12pt]
    Department of Mathematics\\
    Johns Hopkins University
\end{center}

\vspace{0.5cm}

\section*{1. Course Overview}

\textbf{Course:} AS.110.106 Calculus I (Biology and Social Sciences)

\textbf{Primary Text:}
\begin{quote}
Calculus for Biology and Medicine, 4th Edition,
C. Neuhauser and M. Roper, Boston: Pearson Education, January 2018.
\end{quote}

\textbf{ISBN-10:} 0134070046

\textbf{ISBN-13:} 978-0134070049

The course develops the fundamental ideas of single-variable calculus, beginning
with a review of functions and exponential models and progressing through limits,
continuity, differentiation, applications of differentiation, integration, and
applications of the integral.

The official syllabus organizes the course into the following major blocks:

\begin{enumerate}
    \item Review of Basic Properties of Functions
    \item Limits and Continuity
    \item Derivatives
    \item Applications of Differentiation
    \item Integration
    \item Applications of the Integral
\end{enumerate}

The first five major blocks occupy approximately ten weeks, while applications
of the integral occupy one or more additional weeks.

% ============================================================
\section{2. Course Structure}
% ============================================================

\begin{longtable}{p{0.08\textwidth} p{0.27\textwidth} p{0.15\textwidth} p{0.40\textwidth}}
\toprule
\textbf{No.} & \textbf{Major Topic} & \textbf{Approx. Time} &
\textbf{Core Sections} \\
\midrule
1 & Review of Basic Properties of Functions & 2 weeks &
Chapter 1; 2.1; 2.2 \\
2 & Limits and Continuity & 2 weeks &
3.1--3.5, with formal definitions from 3.6 \\
3 & Derivatives & 3 weeks &
4.1--4.11 \\
4 & Applications of Differentiation & 2+ weeks &
5.1--5.5 \\
5 & Integration & 2 weeks &
5.10; 6.1--6.3 \\
6 & Applications of the Integral & 1+ week &
7.1--7.3 \\
\bottomrule
\end{longtable}

\important{The time allocations above follow the official syllabus.}

% ============================================================
\section{3. Prerequisite Knowledge}
% ============================================================

Before beginning the course, students should be comfortable with the following
mathematical ideas. These are prerequisite competencies needed to successfully
engage with the syllabus topics.

\subsection{Algebraic Skills}

Students should be able to:

\begin{itemize}
    \item Simplify algebraic expressions.
    \item Factor polynomials.
    \item Solve linear and quadratic equations.
    \item Solve basic systems of equations.
    \item Work with rational expressions.
    \item Manipulate fractions and exponents.
    \item Use laws of exponents.
    \item Simplify radicals.
    \item Rearrange equations to solve for a specified variable.
\end{itemize}

\subsection{Function Skills}

Students should understand:

\begin{itemize}
    \item Function notation.
    \item Domain and range.
    \item Graphs of functions.
    \item Linear and polynomial functions.
    \item Rational functions.
    \item Exponential functions.
    \item Basic trigonometric functions.
    \item Composition of functions.
    \item Inverse functions.
\end{itemize}

\subsection{Trigonometric Skills}

Students should be familiar with:

\begin{itemize}
    \item Sine, cosine, and tangent.
    \item Basic trigonometric identities.
    \item The unit circle.
    \item Important trigonometric limits and identities.
    \item Graphs of the basic trigonometric functions.
\end{itemize}

% ============================================================
\section{4. Major Course Learning Objectives}
% ============================================================

By the end of the course, students should be able to:

\begin{enumerate}
    \item Analyze and manipulate common classes of functions.
    \item Interpret exponential growth and decay models.
    \item Work with sequences and their limiting behavior.
    \item Understand limits as the foundation of differential and integral calculus.
    \item Determine whether functions are continuous.
    \item Evaluate limits algebraically and graphically.
    \item Understand the formal definition of a limit.
    \item Apply the Sandwich Theorem to appropriate limit problems.
    \item Compute derivatives using the definition of the derivative.
    \item Differentiate functions using standard differentiation rules.
    \item Apply the product, quotient, and chain rules.
    \item Differentiate implicitly defined functions.
    \item Calculate and interpret higher derivatives.
    \item Differentiate trigonometric, exponential, and inverse functions.
    \item Construct linear approximations.
    \item Analyze extrema, monotonicity, concavity, and inflection points.
    \item Apply the Extreme Value and Mean Value Theorems.
    \item Solve optimization problems.
    \item Evaluate appropriate indeterminate limits using L'Hopital's Rule.
    \item Understand antiderivatives and their relationship to derivatives.
    \item Evaluate definite integrals.
    \item Understand and apply the Fundamental Theorem of Calculus.
    \item Interpret integration in applications.
    \item Use the substitution rule.
    \item Apply integration by parts.
    \item Integrate appropriate rational functions using partial fractions.
\end{enumerate}

% ============================================================
\section{5. Module 1: Review of Basic Properties of Functions}
% ============================================================

\textbf{Approximate duration: 2 weeks}

\textbf{Syllabus coverage:} Chapter 1, Section 2.1, Section 2.2.

\subsection{Module Purpose}

This module establishes the algebraic and functional foundation required for
limits, derivatives, and integration.

\subsection{Chapter 1: Basic Properties of Functions}

\subsubsection{Major Concepts}

\begin{itemize}
    \item Definition of a function.
    \item Independent and dependent variables.
    \item Domain and range.
    \item Function notation.
    \item Graphs of functions.
    \item Transformations of functions.
    \item Polynomial functions.
    \item Rational functions.
    \item Exponential functions.
    \item Trigonometric functions.
    \item Composition of functions.
    \item Inverse functions.
\end{itemize}

\subsubsection{Learning Objectives}

Students should be able to:

\begin{enumerate}
    \item Identify the domain and range of a function.
    \item Evaluate functions at specified inputs.
    \item Interpret graphs of functions.
    \item Determine whether a relation represents a function.
    \item Compose two or more functions.
    \item Identify and work with inverse functions.
    \item Analyze common transformations of graphs.
    \item Recognize important classes of functions used in calculus.
\end{enumerate}

\subsection{Section 2.1: Exponential Growth and Decay}

\subsubsection{Major Concepts}

\begin{itemize}
    \item Exponential growth.
    \item Exponential decay.
    \item Growth and decay rates.
    \item Exponential models.
    \item Interpretation of parameters.
    \item Biological and social-science applications.
\end{itemize}

\subsubsection{Learning Objectives}

Students should be able to:

\begin{enumerate}
    \item Recognize exponential growth and decay models.
    \item Interpret the parameters of an exponential model.
    \item Determine whether a model describes growth or decay.
    \item Use exponential functions to model changing quantities.
\end{enumerate}

\subsection{Section 2.2: Sequences}

\subsubsection{Major Concepts}

\begin{itemize}
    \item Definition of a sequence.
    \item Explicit and recursive descriptions.
    \item Sequence notation.
    \item Behavior of sequences.
    \item Convergence and divergence.
\end{itemize}

\subsubsection{Learning Objectives}

Students should be able to:

\begin{enumerate}
    \item Represent sequences using appropriate notation.
    \item Identify patterns in sequences.
    \item Analyze whether sequences converge or diverge.
    \item Connect sequence behavior to the broader concept of limits.
\end{enumerate}

\subsection{Difficult Topics}

\begin{itemize}
    \item Domain restrictions.
    \item Composition and inverse functions.
    \item Exponential growth and decay parameters.
    \item Connecting sequence behavior to limits.
\end{itemize}

% ============================================================
\section{6. Module 2: Limits and Continuity}
% ============================================================

\textbf{Approximate duration: 2 weeks}

\textbf{Syllabus coverage:} Sections 3.1--3.5, with formal definitions
from Section 3.6.

\subsection{Module Purpose}

This module introduces the limiting process that forms the conceptual foundation
for both differentiation and integration.

\subsection{Section 3.1: Limits}

\subsubsection{Major Concepts}

\begin{itemize}
    \item Intuitive interpretation of a limit.
    \item One-sided limits.
    \item Two-sided limits.
    \item Limits from graphs.
    \item Limits from algebraic expressions.
    \item Formal definition of a limit.
\end{itemize}

\subsubsection{Learning Objectives}

Students should be able to:

\begin{enumerate}
    \item Interpret limits graphically.
    \item Evaluate limits algebraically.
    \item Distinguish one-sided from two-sided limits.
    \item Determine when a limit does not exist.
    \item Explain the formal meaning of a limit.
\end{enumerate}

\subsection{Section 3.2: Continuity}

\subsubsection{Major Concepts}

\begin{itemize}
    \item Continuity at a point.
    \item Continuity on an interval.
    \item Discontinuities.
    \item Removable discontinuities.
    \item Infinite discontinuities.
    \item Relationship between continuity and limits.
\end{itemize}

\subsubsection{Learning Objectives}

Students should be able to:

\begin{enumerate}
    \item Determine whether a function is continuous at a point.
    \item Identify discontinuities from graphs and formulas.
    \item Use limits to test continuity.
    \item Explain the relationship between continuity and limiting behavior.
\end{enumerate}

\subsection{Section 3.3: Limits at Infinity}

\subsubsection{Major Concepts}

\begin{itemize}
    \item Limits as $x \to \infty$.
    \item Limits as $x \to -\infty$.
    \item End behavior.
    \item Horizontal asymptotic behavior.
    \item Formal definition of limits at infinity.
\end{itemize}

\subsubsection{Learning Objectives}

Students should be able to:

\begin{enumerate}
    \item Evaluate limits at infinity.
    \item Determine the end behavior of functions.
    \item Interpret horizontal asymptotic behavior.
    \item Relate algebraic structure to behavior at infinity.
\end{enumerate}

\subsection{Section 3.4: Trigonometric Limits and the Sandwich Theorem}

\subsubsection{Major Concepts}

\begin{itemize}
    \item Fundamental trigonometric limits.
    \item Sandwich Theorem.
    \item Bounding functions.
    \item Oscillatory behavior.
\end{itemize}

\subsubsection{Learning Objectives}

Students should be able to:

\begin{enumerate}
    \item Evaluate important trigonometric limits.
    \item Recognize when direct substitution is insufficient.
    \item Construct useful bounds for a function.
    \item Apply the Sandwich Theorem.
\end{enumerate}

\subsection{Section 3.5: Properties of Continuous Functions}

\subsubsection{Major Concepts}

\begin{itemize}
    \item Algebraic properties of continuous functions.
    \item Continuity of sums and products.
    \item Continuity of quotients where defined.
    \item Continuity of compositions.
    \item Consequences of continuity.
\end{itemize}

\subsubsection{Learning Objectives}

Students should be able to:

\begin{enumerate}
    \item Determine continuity using known continuous functions.
    \item Apply properties of continuous functions.
    \item Analyze continuity of composite functions.
\end{enumerate}

\subsection{Difficult Topics}

\begin{itemize}
    \item Formal $\epsilon$--$\delta$ interpretation of limits.
    \item Distinguishing a function value from a limiting value.
    \item One-sided versus two-sided limits.
    \item Limits at infinity.
    \item Trigonometric limits.
    \item Applying the Sandwich Theorem.
\end{itemize}

% ============================================================
\section{7. Module 3: Derivatives}
% ============================================================

\textbf{Approximate duration: 3 weeks}

\textbf{Syllabus coverage:} Sections 4.1--4.11.

\subsection{Module Purpose}

This module develops the derivative as a limit and then builds the principal
rules required for efficient differentiation.

\subsection{Section 4.1: Formal Definition of the Derivative}

\subsubsection{Major Concepts}

\begin{itemize}
    \item Difference quotient.
    \item Derivative as a limit.
    \item Instantaneous rate of change.
    \item Tangent line.
    \item Derivative notation.
\end{itemize}

\subsubsection{Learning Objectives}

Students should be able to:

\begin{enumerate}
    \item Define the derivative using a limit.
    \item Calculate derivatives from first principles.
    \item Interpret the derivative as an instantaneous rate of change.
    \item Interpret the derivative geometrically as a tangent slope.
\end{enumerate}

\subsection{Section 4.2: Basic Rules of Derivatives}

\subsubsection{Major Concepts}

\begin{itemize}
    \item Constant functions.
    \item Basic differentiation rules.
    \item Linearity of differentiation.
\end{itemize}

\subsection{Section 4.3: Power Rule and Rules of Differentiation}

\subsubsection{Major Concepts}

\begin{itemize}
    \item Power Rule.
    \item Constant multiple rule.
    \item Sum and difference rules.
    \item Polynomial differentiation.
\end{itemize}

\subsection{Section 4.4: Product and Quotient Rules}

\subsubsection{Major Concepts}

\begin{itemize}
    \item Product Rule.
    \item Quotient Rule.
    \item Differentiation of rational expressions.
\end{itemize}

\subsection{Section 4.5: Chain Rule}

\subsubsection{Major Concepts}

\begin{itemize}
    \item Composite functions.
    \item Inner and outer functions.
    \item Chain Rule.
    \item Nested differentiation.
\end{itemize}

\subsubsection{Learning Objectives}

Students should be able to:

\begin{enumerate}
    \item Identify composite functions.
    \item Apply the Chain Rule correctly.
    \item Differentiate nested algebraic, trigonometric, and exponential
          expressions.
\end{enumerate}

\subsection{Section 4.6: Implicit Functions and Implicit Differentiation}

\subsubsection{Major Concepts}

\begin{itemize}
    \item Explicit versus implicit functions.
    \item Differentiating equations involving $x$ and $y$.
    \item Solving for $\frac{dy}{dx}$.
\end{itemize}

\subsubsection{Learning Objectives}

Students should be able to:

\begin{enumerate}
    \item Differentiate implicitly defined relationships.
    \item Solve for $\frac{dy}{dx}$.
    \item Interpret implicit derivatives.
\end{enumerate}

\subsection{Section 4.7: Higher Derivatives}

\subsubsection{Major Concepts}

\begin{itemize}
    \item Second derivative.
    \item Higher-order derivatives.
    \item Notation for higher derivatives.
    \item Interpretation of derivatives of derivatives.
\end{itemize}

\subsection{Section 4.8: Derivatives of Trigonometric Functions}

\subsubsection{Major Concepts}

\begin{itemize}
    \item Derivatives of sine and cosine.
    \item Derivatives of other basic trigonometric functions.
    \item Combining trigonometric derivatives with differentiation rules.
\end{itemize}

\subsection{Section 4.9: Derivatives of Exponential Functions}

\subsubsection{Major Concepts}

\begin{itemize}
    \item Derivative of exponential functions.
    \item The role of $e$.
    \item Exponential growth rates.
    \item Chain Rule with exponential functions.
\end{itemize}

\subsection{Section 4.10: Derivatives of Inverse Functions}

\subsubsection{Major Concepts}

\begin{itemize}
    \item Inverse functions.
    \item Derivatives of inverse functions.
    \item Relationship between inverse functions and their derivatives.
\end{itemize}

\subsection{Section 4.11: Linear Approximation}

\subsubsection{Major Concepts}

\begin{itemize}
    \item Tangent-line approximation.
    \item Linearization.
    \item Approximation of function values.
    \item Error interpretation.
\end{itemize}

\subsubsection{Learning Objectives}

Students should be able to:

\begin{enumerate}
    \item Construct a linear approximation.
    \item Use linearization to estimate function values.
    \item Determine when linear approximation is appropriate.
\end{enumerate}

\subsection{Difficult Topics}

The principal conceptual and computational challenges in this module are:

\begin{itemize}
    \item Understanding the derivative as a limit.
    \item Moving from first-principles differentiation to derivative rules.
    \item Correct application of the Chain Rule.
    \item Combining multiple differentiation rules.
    \item Implicit differentiation.
    \item Derivatives of inverse functions.
    \item Higher derivatives.
    \item Choosing an appropriate linear approximation.
\end{itemize}

% ============================================================
\section{8. Module 4: Applications of Differentiation}
% ============================================================

\textbf{Approximate duration: 2+ weeks}

\textbf{Syllabus coverage:} Sections 5.1--5.5.

\subsection{Module Purpose}

This module moves from computing derivatives to using derivatives to analyze
and solve mathematical problems.

\subsection{Section 5.1: Extreme and Mean Value Theorems}

\subsubsection{Major Concepts}

\begin{itemize}
    \item Extreme Value Theorem.
    \item Mean Value Theorem.
    \item Conditions required by the theorems.
    \item Interpretation of theorem conclusions.
\end{itemize}

\subsubsection{Learning Objectives}

Students should be able to:

\begin{enumerate}
    \item State the hypotheses of the Extreme Value Theorem.
    \item State the hypotheses of the Mean Value Theorem.
    \item Determine whether the hypotheses are satisfied.
    \item Interpret the conclusions geometrically.
\end{enumerate}

\subsection{Section 5.2: Monotonicity and Concavity}

\subsubsection{Major Concepts}

\begin{itemize}
    \item Increasing functions.
    \item Decreasing functions.
    \item First derivative test.
    \item Concavity.
    \item Second derivative.
    \item Relationship between derivatives and graph shape.
\end{itemize}

\subsubsection{Learning Objectives}

Students should be able to:

\begin{enumerate}
    \item Determine intervals of increase and decrease.
    \item Use the first derivative to analyze a function.
    \item Determine intervals of concavity.
    \item Use the second derivative to analyze graph shape.
\end{enumerate}

\subsection{Section 5.3: Extrema and Inflection Points}

\subsubsection{Major Concepts}

\begin{itemize}
    \item Local extrema.
    \item Absolute extrema.
    \item Critical points.
    \item Inflection points.
    \item First and second derivative tests.
\end{itemize}

\subsubsection{Learning Objectives}

Students should be able to:

\begin{enumerate}
    \item Locate critical points.
    \item Classify local extrema.
    \item Determine absolute extrema when appropriate.
    \item Identify inflection points.
    \item Construct qualitative graphs using derivative information.
\end{enumerate}

\subsection{Section 5.4: Optimization}

\subsubsection{Major Concepts}

\begin{itemize}
    \item Mathematical formulation of optimization problems.
    \item Objective functions.
    \item Constraints.
    \item Critical points.
    \item Maximum and minimum values.
    \item Interpretation of solutions.
\end{itemize}

\subsubsection{Learning Objectives}

Students should be able to:

\begin{enumerate}
    \item Translate a verbal problem into a mathematical model.
    \item Identify the quantity to maximize or minimize.
    \item Express the objective in terms of one variable when appropriate.
    \item Use derivatives to locate optimal values.
    \item Interpret the result in the original context.
\end{enumerate}

\subsection{Section 5.5: L'Hopital's Rule}

\subsubsection{Major Concepts}

\begin{itemize}
    \item Indeterminate forms.
    \item Limits involving quotients.
    \item L'Hopital's Rule.
    \item Repeated application of L'Hopital's Rule.
\end{itemize}

\subsubsection{Learning Objectives}

Students should be able to:

\begin{enumerate}
    \item Recognize appropriate indeterminate forms.
    \item Determine when L'Hopital's Rule can be applied.
    \item Evaluate appropriate limits using the rule.
    \item Recognize when an alternative limit technique is preferable.
\end{enumerate}

\subsection{Difficult Topics}

\begin{itemize}
    \item Applying theorem hypotheses correctly.
    \item Distinguishing local and absolute extrema.
    \item Interpreting first and second derivative information.
    \item Identifying inflection points.
    \item Translating word problems into optimization models.
    \item Correct use of L'Hopital's Rule.
\end{itemize}

% ============================================================
\section{9. Module 5: Integration}
% ============================================================

\textbf{Approximate duration: 2 weeks}

\textbf{Syllabus coverage:} Section 5.10 and Sections 6.1--6.3.

\subsection{Module Purpose}

This module introduces integration as the inverse process of differentiation
and as a method for accumulating quantities.

\subsection{Section 5.10: Antiderivatives}

\subsubsection{Major Concepts}

\begin{itemize}
    \item Antiderivatives.
    \item Families of antiderivatives.
    \item Indefinite integrals.
    \item Constant of integration.
    \item Relationship between differentiation and integration.
\end{itemize}

\subsubsection{Learning Objectives}

Students should be able to:

\begin{enumerate}
    \item Identify an antiderivative of a function.
    \item Compute basic indefinite integrals.
    \item Include the constant of integration.
    \item Verify an antiderivative by differentiation.
\end{enumerate}

\subsection{Section 6.1: The Definite Integral}

\subsubsection{Major Concepts}

\begin{itemize}
    \item Definite integral.
    \item Accumulation.
    \item Signed area.
    \item Riemann-sum interpretation.
    \item Integral notation.
\end{itemize}

\subsubsection{Learning Objectives}

Students should be able to:

\begin{enumerate}
    \item Interpret the definite integral as accumulated quantity.
    \item Interpret the definite integral geometrically.
    \item Distinguish signed area from total geometric area.
    \item Evaluate basic definite integrals.
\end{enumerate}

\subsection{Section 6.2: Fundamental Theorem of Calculus}

\subsubsection{Major Concepts}

\begin{itemize}
    \item Relationship between derivatives and integrals.
    \item Fundamental Theorem of Calculus.
    \item Evaluating definite integrals using antiderivatives.
    \item Accumulation functions.
\end{itemize}

\subsubsection{Learning Objectives}

Students should be able to:

\begin{enumerate}
    \item Explain the relationship between differentiation and integration.
    \item Apply the Fundamental Theorem of Calculus.
    \item Evaluate definite integrals using antiderivatives.
    \item Differentiate accumulation functions when appropriate.
\end{enumerate}

\subsection{Section 6.3: Applications of Integration}

\subsubsection{Major Concepts}

\begin{itemize}
    \item Accumulated change.
    \item Area.
    \item Total quantities.
    \item Interpretation of integrals in applications.
\end{itemize}

\subsubsection{Learning Objectives}

Students should be able to:

\begin{enumerate}
    \item Formulate integrals representing accumulated quantities.
    \item Interpret the result of an integral in context.
    \item Use integration to solve basic applied problems.
\end{enumerate}

\subsection{Difficult Topics}

\begin{itemize}
    \item Understanding the definite integral conceptually.
    \item Distinguishing indefinite and definite integrals.
    \item Understanding signed area.
    \item Understanding the Fundamental Theorem of Calculus.
    \item Translating an accumulation problem into an integral.
\end{itemize}

% ============================================================
\section{10. Module 6: Applications of the Integral}
% ============================================================

\textbf{Approximate duration: 1+ week}

\textbf{Syllabus coverage:} Sections 7.1--7.3.

\subsection{Module Purpose}

This module develops techniques for evaluating integrals that cannot be handled
directly by the most basic antiderivative rules.

\subsection{Section 7.1: Substitution Rule}

\subsubsection{Major Concepts}

\begin{itemize}
    \item Change of variables.
    \item $u$-substitution.
    \item Reverse Chain Rule.
    \item Definite integrals under substitution.
\end{itemize}

\subsubsection{Learning Objectives}

Students should be able to:

\begin{enumerate}
    \item Recognize integrals suitable for substitution.
    \item Select an appropriate substitution.
    \item Transform an integral using the substitution rule.
    \item Evaluate definite integrals using substitution.
\end{enumerate}

\subsection{Section 7.2: Integration by Parts and Practicing Integration}

\subsubsection{Major Concepts}

\begin{itemize}
    \item Integration by parts.
    \item Product Rule as the basis for integration by parts.
    \item Choosing $u$ and $dv$.
    \item Combining integration techniques.
    \item Strategic selection of an integration method.
\end{itemize}

\subsubsection{Learning Objectives}

Students should be able to:

\begin{enumerate}
    \item Apply the integration-by-parts formula.
    \item Choose appropriate factors for $u$ and $dv$.
    \item Recognize when integration by parts is useful.
    \item Combine integration methods when necessary.
\end{enumerate}

\subsection{Section 7.3: Rational Functions and Partial Fractions}

\subsubsection{Major Concepts}

\begin{itemize}
    \item Rational functions.
    \item Partial fraction decomposition.
    \item Decomposition of rational expressions.
    \item Integration using partial fractions.
\end{itemize}

\subsubsection{Learning Objectives}

Students should be able to:

\begin{enumerate}
    \item Recognize when partial fractions are appropriate.
    \item Decompose suitable rational functions.
    \item Integrate the resulting simpler expressions.
\end{enumerate}

\subsection{Difficult Topics}

\begin{itemize}
    \item Recognizing which integration technique to use.
    \item Correctly performing substitutions.
    \item Integration by parts and choosing $u$ and $dv$.
    \item Partial fraction decomposition.
    \item Combining multiple integration techniques.
\end{itemize}

% ============================================================
\section{11. Conceptual Dependency Structure}
% ============================================================

The major conceptual dependencies of the course can be summarized as follows:

\[
\boxed{
\text{Functions}
\longrightarrow
\text{Limits}
\longrightarrow
\text{Continuity}
\longrightarrow
\text{Derivatives}
}
\]

\[
\boxed{
\text{Derivatives}
\longrightarrow
\text{Derivative Applications}
\longrightarrow
\text{Optimization and Graph Analysis}
}
\]

\[
\boxed{
\text{Derivatives}
\longrightarrow
\text{Antiderivatives}
\longrightarrow
\text{Definite Integrals}
\longrightarrow
\text{Fundamental Theorem of Calculus}
}
\]

\[
\boxed{
\text{Integration}
\longrightarrow
\text{Substitution}
\longrightarrow
\text{Integration by Parts}
\longrightarrow
\text{Partial Fractions}
}
\]

The overall conceptual progression is therefore:

\[
\text{Algebra and Functions}
\rightarrow
\text{Limits}
\rightarrow
\text{Continuity}
\rightarrow
\text{Derivatives}
\rightarrow
\text{Applications}
\rightarrow
\text{Integration}
\rightarrow
\text{Integration Techniques}.
\]

% ============================================================
\section{12. Difficult Topics Across the Entire Course}
% ============================================================

The following topics are likely to require the greatest conceptual effort
because they introduce new mathematical reasoning or combine several earlier
ideas.

\begin{enumerate}
    \item \textbf{Formal definition of the limit}

    Students must move beyond numerical or graphical intuition toward a precise
    mathematical description of limiting behavior.

    \item \textbf{Trigonometric limits and the Sandwich Theorem}

    These require students to recognize appropriate transformations and
    bounding arguments.

    \item \textbf{Derivative from first principles}

    Students must connect the difference quotient to the geometric and
    instantaneous-rate interpretations of the derivative.

    \item \textbf{Chain Rule}

    Correctly identifying nested functions and differentiating from the
    outside inward is a major transition from basic derivative rules.

    \item \textbf{Implicit differentiation}

    Students must differentiate equations in which the dependent variable is
    not explicitly isolated.

    \item \textbf{Applications of derivatives}

    These require interpretation rather than simply executing a derivative
    calculation.

    \item \textbf{Optimization}

    Optimization requires translating a real-world problem into a mathematical
    model before calculus can be applied.

    \item \textbf{L'Hopital's Rule}

    Students must recognize appropriate indeterminate forms and understand
    when the theorem is applicable.

    \item \textbf{Definite integral and the Fundamental Theorem of Calculus}

    This is the central conceptual transition from differentiation to
    accumulation.

    \item \textbf{Choosing integration techniques}

    The challenge is often not performing an integration technique but
    recognizing which technique is appropriate.

    \item \textbf{Partial fractions}

    This combines algebraic decomposition with integration.
\end{enumerate}

% ============================================================
\section{13. Course-Wide Skill Progression}
% ============================================================

\begin{longtable}{p{0.20\textwidth} p{0.35\textwidth} p{0.35\textwidth}}
\toprule
\textbf{Stage} & \textbf{Primary Skill} & \textbf{Expected Student Capability} \\
\midrule
Foundation &
Function analysis &
Manipulate, graph, compose, and interpret functions \\

Limits &
Limiting reasoning &
Evaluate and interpret limiting behavior \\

Continuity &
Local behavior &
Determine whether functions are continuous \\

Differentiation &
Rate of change &
Compute and interpret derivatives \\

Derivative applications &
Analysis &
Use derivatives to study graphs and solve optimization problems \\

Integration &
Accumulation &
Compute and interpret definite and indefinite integrals \\

Fundamental Theorem &
Unification &
Connect differentiation and integration \\

Integration techniques &
Strategic problem solving &
Select and apply appropriate integration methods \\
\bottomrule
\end{longtable}

% ============================================================
\section{14. Suggested Module-by-Module Master Checklist}
% ============================================================

\subsection*{Module 1 -- Functions}

\begin{itemize}[label=$\square$]
    \item Function notation
    \item Domain and range
    \item Graphs
    \item Function transformations
    \item Polynomial and rational functions
    \item Exponential functions
    \item Trigonometric functions
    \item Composition
    \item Inverse functions
    \item Exponential growth and decay
    \item Sequences
\end{itemize}

\subsection*{Module 2 -- Limits and Continuity}

\begin{itemize}[label=$\square$]
    \item One-sided limits
    \item Two-sided limits
    \item Algebraic limits
    \item Graphical limits
    \item Formal definition of limit
    \item Continuity
    \item Discontinuities
    \item Limits at infinity
    \item Trigonometric limits
    \item Sandwich Theorem
    \item Properties of continuous functions
\end{itemize}

\subsection*{Module 3 -- Derivatives}

\begin{itemize}[label=$\square$]
    \item Definition of derivative
    \item Difference quotient
    \item Tangent lines
    \item Basic derivative rules
    \item Power Rule
    \item Product Rule
    \item Quotient Rule
    \item Chain Rule
    \item Implicit differentiation
    \item Higher derivatives
    \item Trigonometric derivatives
    \item Exponential derivatives
    \item Inverse-function derivatives
    \item Linear approximation
\end{itemize}

\subsection*{Module 4 -- Applications of Differentiation}

\begin{itemize}[label=$\square$]
    \item Extreme Value Theorem
    \item Mean Value Theorem
    \item Increasing/decreasing functions
    \item Monotonicity
    \item Concavity
    \item Critical points
    \item Local extrema
    \item Absolute extrema
    \item Inflection points
    \item Curve analysis
    \item Optimization
    \item L'Hopital's Rule
\end{itemize}

\subsection*{Module 5 -- Integration}

\begin{itemize}[label=$\square$]
    \item Antiderivatives
    \item Indefinite integrals
    \item Definite integrals
    \item Accumulation
    \item Signed area
    \item Riemann-sum interpretation
    \item Fundamental Theorem of Calculus
    \item Applications of integration
\end{itemize}

\subsection*{Module 6 -- Applications of the Integral}

\begin{itemize}[label=$\square$]
    \item Substitution Rule
    \item Change of variables
    \item Integration by parts
    \item Choosing integration methods
    \item Rational functions
    \item Partial fractions
\end{itemize}

% ============================================================
\section{15. Final Course Learning Outcomes}
% ============================================================

Upon successful completion of AS.110.106 Calculus I, a student should be able
to demonstrate proficiency in the following areas:

\begin{enumerate}
    \item \textbf{Function Analysis:}
    Analyze, manipulate, and interpret the principal classes of functions
    used throughout introductory calculus.

    \item \textbf{Limit Reasoning:}
    Evaluate and interpret limits and use limits to understand continuity and
    the derivative.

    \item \textbf{Continuity:}
    Determine and justify continuity and identify different types of
    discontinuous behavior.

    \item \textbf{Differentiation:}
    Compute derivatives using the definition of the derivative and the
    standard differentiation rules.

    \item \textbf{Derivative Interpretation:}
    Interpret derivatives geometrically and as rates of change.

    \item \textbf{Function Analysis Using Derivatives:}
    Use derivatives to determine monotonicity, concavity, extrema, and
    inflection points.

    \item \textbf{Optimization:}
    Formulate and solve optimization problems arising from mathematical and
    applied contexts.

    \item \textbf{Integration:}
    Compute antiderivatives and definite integrals and interpret integration
    as accumulation.

    \item \textbf{Fundamental Theorem:}
    Explain and apply the Fundamental Theorem of Calculus to connect
    differentiation and integration.

    \item \textbf{Applied Integration:}
    Formulate and interpret integrals representing accumulated quantities.

    \item \textbf{Integration Techniques:}
    Apply substitution, integration by parts, and partial fractions to
    evaluate appropriate integrals.

    \item \textbf{Mathematical Modeling:}
    Translate biological, social-science, and mathematical situations into
    appropriate functional, differential, or integral models.
\end{enumerate}

% ============================================================
\section{16. Complete Course Map}
% ============================================================

\begin{center}
\begin{tabular}{c}
\textbf{MODULE 1: FUNCTIONS}\\
\text{Functions $\rightarrow$ Exponential Growth/Decay $\rightarrow$ Sequences}\\
$\downarrow$\\[4pt]
\textbf{MODULE 2: LIMITS AND CONTINUITY}\\
\text{Limits $\rightarrow$ Continuity $\rightarrow$ Limits at Infinity}\\
\text{Trigonometric Limits $\rightarrow$ Sandwich Theorem}\\
$\downarrow$\\[4pt]
\textbf{MODULE 3: DERIVATIVES}\\
\text{Definition $\rightarrow$ Rules $\rightarrow$ Product/Quotient}\\
\text{Chain Rule $\rightarrow$ Implicit Differentiation}\\
\text{Higher Derivatives $\rightarrow$ Special Functions $\rightarrow$ Linear Approximation}\\
$\downarrow$\\[4pt]
\textbf{MODULE 4: APPLICATIONS OF DIFFERENTIATION}\\
\text{EVT/MVT $\rightarrow$ Monotonicity $\rightarrow$ Concavity}\\
\text{Extrema $\rightarrow$ Inflection Points $\rightarrow$ Optimization}\\
\text{L'Hopital's Rule}\\
$\downarrow$\\[4pt]
\textbf{MODULE 5: INTEGRATION}\\
\text{Antiderivatives $\rightarrow$ Definite Integral}\\
\text{$\rightarrow$ Fundamental Theorem of Calculus}\\
\text{$\rightarrow$ Applications of Integration}\\
$\downarrow$\\[4pt]
\textbf{MODULE 6: APPLICATIONS OF THE INTEGRAL}\\
\text{Substitution $\rightarrow$ Integration by Parts}\\
\text{$\rightarrow$ Rational Functions $\rightarrow$ Partial Fractions}
\end{tabular}
\end{center}

\vfill

\begin{center}
\textit{End of Course Structure}
\end{center}

\end{document}